\chapter{Plant, Machinery and Vehicle Risk Assessment}
\label{chap:Plant, Machinery and Vehicle Risk Assessment}

%---------------------------- SECTION ----------------------------
\section{Why this assessment matters in construction}

Being struck by a moving vehicle or piece of construction plant is the second most common cause of fatal injury in the UK construction industry, responsible for approximately 20 per cent of all construction deaths in any given year. The HSE's analysis of plant-related fatalities consistently identifies a predictable pattern: a pedestrian worker in a location that the plant operator cannot see, a construction vehicle reversing without adequate supervision, a pedestrian route that was not separated from the plant operating area, and a risk assessment that acknowledged the vehicle/pedestrian conflict without implementing an effective management response. The predictability of this pattern is both its most sobering and its most actionable characteristic: these deaths are not random events; they are the consequence of known hazards meeting inadequate controls.

The variety of construction plant operating on a typical site creates multiple simultaneous exposure points. Mobile cranes, tower cranes, telehandlers, forklifts, mechanical excavators, dumpers, road rollers, articulated trucks, concrete pump trucks, and site welfare vehicles may all be present on a busy site at the same time. Each has its own swept path, its own blind spots, its own operational envelope, and its own proximity risk to pedestrian workers. The risk assessment must address each type of plant individually and must address the cumulative effect of multiple plant types operating simultaneously in proximity to pedestrian routes.

The LOLER 1998 thorough examination regime for lifting plant, the PUWER 1998 inspection requirements for all work equipment, and the CPCS (Construction Plant Competence Scheme) operator certification system collectively form the framework within which plant risk management operates. But legislation and certification cannot substitute for the fundamental control of separating pedestrians from moving plant: this is an engineering and site organisation measure that must be planned, implemented, and maintained as a primary control, not a desirable enhancement.

The management of reversing vehicles deserves specific attention because reversing operations account for a disproportionate share of struck-by fatalities. Site layouts that require vehicles to reverse over long distances, that place welfare facilities or pedestrian access points in reversing paths, or that do not provide adequate banksman supervision create conditions in which a fatal reversing incident becomes statistically probable given sufficient vehicle movements and time. The risk assessment must address reversing management as a specific topic.

\barquote{``Separation of pedestrians from plant is an engineering control, not a guidance document. A barrier prevents a collision; a briefing note does not.''}

\example{Site Reality}{A civil engineering contractor is constructing a new access road on a commercial development. Three 6-tonne site dumpers are in continuous operation between the spoil excavation area and the tip, running along a haul route that crosses the main pedestrian access path from the site entrance to the welfare facilities. The site manager has written ``pedestrians to give way to plant'' on the site induction document. No physical segregation is installed at the crossing point. A labourer crosses the haul route on his way to the welfare cabin; a dumper reversing out of the tip area strikes him from behind, causing fatal crush injuries. The investigation identifies the absence of physical segregation at the crossing point, the absence of banksman supervision, and a site layout that placed the welfare facilities on the opposite side of the haul route from the site entrance as primary causal factors. The risk assessment identified the pedestrian/plant crossing as a hazard; the control response was wholly inadequate.}

%---------------------------- SECTION ----------------------------
\section{Legal and professional framework}

The principal legislative framework for plant, machinery and vehicle safety in construction comprises PUWER 1998 (general work equipment), LOLER 1998 (lifting equipment), the Health and Safety at Work etc.\ Act 1974, and CDM 2015. The CDM 2015 requires the Construction Phase Plan to address the arrangements for managing plant and vehicle movements, including the segregation of pedestrians and vehicles, the management of reversing operations, and the organisation of delivery and haul routes. HSE guidance HSG144 (Safe Use of Vehicles on Construction Sites) is the primary reference document for vehicle safety management in construction.

The Provision and Use of Work Equipment Regulations 1998 require that all plant and machinery is suitable for the task, maintained in an efficient state of repair, inspected at appropriate intervals, and operated only by persons who have received adequate training and are competent to do so. LOLER 1998 requires thorough examination of all lifting plant (cranes, telehandlers in lifting mode, excavators used for lifting, forklifts) by a competent person at prescribed intervals (six-monthly for equipment used to lift persons; annually for other lifting equipment) and that all lifting operations are properly planned by a competent person.

\paragraph{Key frameworks relevant to this chapter include: PUWER 1998; LOLER 1998; CDM 2015; HSG144; CPCS competency scheme.}

\begin{itemize}
  \bitem{PUWER 1998}{Requires that all construction plant and vehicles are suitable for their purpose, maintained in efficient working order, inspected at appropriate intervals, and operated only by persons with adequate training and competency; applies to mechanical excavators, dumpers, rollers, forklifts, telehandlers, and all mobile plant.}
  \bitem{LOLER 1998}{Requires thorough examination of all lifting plant by a competent person at prescribed intervals; requires that all lifting operations are properly planned, appropriately supervised, and carried out safely; the thorough examination certificate must be on site and in date before any lifting operation commences.}
  \bitem{CDM 2015}{Requires the Construction Phase Plan to address vehicle and plant movement management, pedestrian/vehicle segregation, banksman arrangements, and delivery logistics; the principal contractor has primary responsibility for coordinating plant movements on the common parts of the site.}
  \bitem{HSG144 Safe Use of Vehicles on Construction Sites}{Provides detailed guidance on vehicle route planning, segregation measures, reversing management, banksman duties and competency, delivery vehicle management, and inspection regimes; compliance with HSG144 represents the minimum acceptable standard for vehicle safety management on UK construction sites.}
\end{itemize}

%---------------------------- SECTION ----------------------------
\section{Conducting the assessment using the five-step method}

\subsection{Step 1 --- Identify hazards}
\paragraph{Plant, machinery and vehicle hazards in construction include: pedestrians struck by moving plant or vehicles in areas where pedestrian and plant routes overlap; pedestrians struck by reversing vehicles in the absence of banksman supervision or proximity warning devices; operatives struck by falling loads during lifting operations where loads are suspended over occupied areas; plant overturning on inadequate or unstable ground, causing crush injury to the operator and persons adjacent; plant operators struck by overhead electricity lines; operatives struck by plant swinging radius (particularly excavator counterweights); mechanical failure of plant causing injury (brake failure, hydraulic failure, tyre blowout); persons trapped between plant and fixed structures or other plant; and vehicle loads improperly secured causing load shedding onto pedestrian areas.}

\subsection{Step 2 --- Decide who or what is affected}
\paragraph{Persons at risk from plant and vehicle hazards include: site operatives working in proximity to plant operating areas; banksmen and lift supervisors working adjacent to plant; delivery drivers whose vehicles are being managed on the site; pedestrians using routes adjacent to haul roads; members of the public adjacent to demolition sites or sites with plant operating adjacent to the site boundary; plant operators themselves (overturning, overhead cables); and emergency services personnel attending plant incidents. Special risk arises for workers with hearing impairment who may not hear audible proximity warnings, and for workers whose hi-visibility clothing has become dirty and is no longer conspicuous.}

\subsection{Step 3 --- Evaluate likelihood and consequence}
\paragraph{Struck-by plant incidents carry fatality as the realistic consequence at the upper end of the spectrum, placing the inherent risk in the Very High category (20--25 on a 5×5 matrix) where pedestrian/plant separation is absent. Even where segregation is partially implemented, the likelihood of incidental proximity between plant and pedestrians on a working construction site is High. Effective physical segregation can reduce this inherent risk dramatically: a site that physically routes all pedestrians through barriers that prevent them from entering plant operating areas, and physically routes plant away from pedestrian areas, can achieve a Low residual risk (4--6/25) for routine operations.}

\subsection{Step 4 --- Select and implement controls}
\paragraph{The primary control for plant and vehicle safety is physical separation of pedestrians from plant: design the site layout to provide distinct, non-overlapping routes for plant and pedestrians; install physical barriers (Heras fencing, Jersey barriers, Armco) to separate plant haul roads from pedestrian routes; create a single controlled crossing point between pedestrian and plant zones managed by a banksman; design the delivery area so that delivery vehicles do not reverse through pedestrian routes to reach their unloading point. Secondary controls include: competency requirements (CPCS card for all plant operators); pre-start plant inspections and maintenance records; banksman requirements for all reversing operations; proximity warning devices on plant; overhead electricity line risk assessment and goal-post systems where lines are adjacent to plant routes.}

\subsection{Step 5 --- Record and review}
\paragraph{The plant and vehicle risk assessment must be reviewed whenever the site layout changes, whenever a new type of plant is introduced, and whenever a new delivery or haulage contractor commences operations. The daily pre-start plant inspection records must be retained on site and reviewed by the site manager weekly. Any near-miss involving plant and pedestrians must trigger an immediate review of the segregation measures and must be reported to the principal contractor before operations resume. LOLER thorough examination certificates must be checked before any new piece of lifting plant commences operations on site.}

\example{Five-Step Application}{A roofing contractor is using a telehandler to lift roofing materials from a ground-level materials area to a first-floor roof deck. Step 1 identifies: load suspended over the pedestrian access route to the first-floor staircase; potential for pedestrian to walk under the load during the lift; telehandler on uneven ground adjacent to the building. Step 2 identifies: two roofers on the roof deck receiving the load; site operatives using the staircase access route; delivery driver in cab of an adjacent vehicle. Step 3 rates inherent struck-by-suspended-load risk as Very High (25/25). Step 4 implements: physical exclusion zone established beneath the lift area using barriers before any lift commences; lift plan prepared and signed by competent person; ground stability check completed before positioning of telehandler; LOLER certificate confirmed in date; two-way radio communication between telehandler operator and roof-deck receiver established; no lift permitted while any person is within the exclusion zone. Residual risk rated Low (6/25). Step 5: lift plan reviewed before each new lift sequence commences.}

%---------------------------- SECTION ----------------------------
\section{Applying the hierarchy of controls}

\paragraph{The hierarchy for plant and vehicle safety must place physical separation as the primary engineering control. Administrative controls and PPE (hi-visibility clothing) are secondary measures that reduce the consequence of an encounter between a person and a plant that should not have been possible in the first place.}

\begin{itemize}
  \bitem{Elimination}{Design the site logistics so that pedestrians and plant do not need to occupy the same space at the same time: use remote material storage areas with crane or telehandler access to avoid plant movement through the working area; use offsite prefabrication delivered as complete assemblies to reduce delivery frequency.}
  \bitem{Substitution}{Replace vehicle reversing with forward-travel-only routes using one-way systems; use dumpers with reversing cameras rather than relying on banksman supervision alone; replace articulated delivery vehicles with smaller vehicles with better visibility where site dimensions permit.}
  \bitem{Engineering controls}{Physical pedestrian/plant segregation using robust barriers (Heras fencing, concrete blocks, Armco crash barriers) along the full length of haul routes; designated and signed pedestrian crossing points managed by a banksman; exclusion zones around plant operating areas; overhead electricity line goal-post systems; proximity alert systems on plant.}
  \bitem{Administrative controls}{CPCS competency requirements for all plant operators; pre-start plant inspection routine with documented records; lift plans for all lifting operations prepared by a competent person; banksman requirements and competency standards for all reversing operations; site induction covering plant movement routes and pedestrian exclusion zones; daily briefing of plant operators on site layout and pedestrian restrictions.}
  \bitem{PPE}{High-visibility vest (EN ISO 20471 Class 2 minimum) for all persons on site as a baseline; hard hat; hearing protection in areas of high plant noise; safety boots. PPE is a last-resort supplementary control, not a substitute for physical segregation.}
\end{itemize}

%---------------------------- SECTION ----------------------------
\section{Cadence of assessment and review triggers}

\paragraph{Plant and vehicle safety assessments must be reviewed as the site programme progresses and the plant fleet, delivery logistics, and site layout evolve. A static assessment prepared at the start of the project will not reflect the conditions that exist once multiple trades and multiple delivery contractors are on site simultaneously.}

\begin{itemize}
  \bitem{Before new plant is introduced}{Any new type of plant or vehicle commencing operations on site must be assessed before first use; its operational envelope, blind spots, reversing requirements, and LOLER status must be confirmed and the existing site segregation measures reviewed for adequacy.}
  \bitem{When site layout changes}{Any significant change in site layout, including the repositioning of welfare facilities, material storage areas, access gates, or haul routes, must trigger a review of pedestrian/plant segregation arrangements.}
  \bitem{At delivery of major loads}{Complex lifting and delivery operations (crane erection, structural steel deliveries, prefabricated module installation) must be individually assessed before the operation; the lift plan and delivery management plan must be prepared in advance by a competent person.}
  \bitem{Following any near-miss or incident}{Any incident involving plant and a person, or any near-miss involving an unplanned proximity between plant and a pedestrian, must trigger an immediate review of the segregation arrangements and must be notified to the principal contractor.}
  \bitem{Weekly site safety inspection}{Plant operating areas, segregation barriers, crossing points, and LOLER documentation must be checked at the weekly site safety inspection; any deficiencies must be allocated for remedy with a defined deadline.}
\end{itemize}

%---------------------------- SECTION ----------------------------
\section{Common failure modes}

\begin{itemize}
  \bitem{Segregation planned but not implemented}{The site layout drawing shows separate plant and pedestrian routes, but the physical barriers are not installed because they have not been procured, or because they have been removed to allow a delivery and not reinstated. The assessment assumes segregation exists; the site conditions do not provide it.}
  \bitem{Banksman not trained or not present}{Banksman supervision for reversing operations is specified in every relevant RAMS, but trained banksmen are not present during all reversing operations because the site staffing model does not provide dedicated banksman resource. Untrained operatives are asked to guide plant; they do not know the correct signals, and the plant operator does not know how to interpret them.}
  \bitem{LOLER certificates expired or not on site}{Lifting plant arrives on site without a current LOLER thorough examination certificate. The site manager does not request the certificate before operations commence. The plant operates for weeks without verified inspection status.}
  \bitem{Pedestrians using plant routes}{Welfare facilities are sited on one side of a haul road; operatives cross the haul road informally throughout the day because there is no designated crossing point and the journey via the formal route would add several minutes. The informal crossing becomes a routine that everyone on site knows about except the site manager.}
  \bitem{Delivery vehicles not managed}{Third-party delivery vehicles arrive at site without specific induction to the site segregation rules; they reverse to unloading points through pedestrian areas using the cab mirrors alone, with no banksman assistance and no exclusion zone management.}
  \bitem{Ground conditions for plant not assessed}{Mobile cranes and telehandlers are positioned without ground-bearing capacity assessment; the ground beneath outrigger pads or wheels is made ground, backfill over drainage runs, or soft clay; the plant overturns or sinks, causing injury to the operator and adjacent persons.}
\end{itemize}

\example{Failure Pattern}{A principal contractor is managing a demolition and new-build project in a town centre. The site access is a single 6-metre gate used by both pedestrians entering the site and demolition vehicles leaving with loaded skips. The site manager instructs operatives to ``be aware of vehicles at the gate'' in the site induction. The banksman resource consists of the site manager's own presence at the gate ``when available''. During week three, a skip lorry exits the gate and strikes a pedestrian operative who was walking in from the car park. The investigation identifies the absence of a pedestrian/vehicle segregated access arrangement as the primary causal factor. The risk assessment had identified the gate as a pedestrian/vehicle conflict point; the control response was an instruction to operatives to be vigilant.}

%---------------------------- CHAPTER SUMMARY ----------------------------
\section{Chapter Summary}

\begin{table}[H]
\centering
\begin{tabularx}{\textwidth}{>{\raggedright\arraybackslash}p{3.2cm}X}
\toprule
\textbf{Assessment Element} & \textbf{Operational Requirement} \\
\midrule
Legal/Professional Basis & PUWER 1998; LOLER 1998; CDM 2015; HSG144. Physical separation of pedestrians from plant is the primary control; LOLER certificates must be verified before lifting plant commences operations. \\
Method & Apply the full five-step process and document inherent/residual risk. \\
Controls & Physical segregation of pedestrian and plant routes; designated managed crossing points; CPCS competency for all plant operators; banksman supervision for reversing; lift plans for lifting operations; pre-start plant inspection records. \\
Cadence & Before new plant is introduced; when site layout changes; before major lifts; following any near-miss; weekly safety inspection. \\
Assurance & Use audit, incident learning, and governance reporting to verify effectiveness. \\
\bottomrule
\end{tabularx}
\end{table}

\begin{itemize}
  \bitem{Key takeaway 1}{Physical separation of pedestrians from plant is an engineering control that must be designed into the site layout before work commences; instructions to pedestrians to ``watch out for plant'' are not an effective control.}
  \bitem{Key takeaway 2}{LOLER thorough examination certificates for all lifting plant must be verified before first use on site; lifting plant without a current certificate must not be operated.}
  \bitem{Key takeaway 3}{Banksman supervision for all reversing operations must be resourced as a planned activity, not delegated informally to the nearest available operative; banksmen must have received formal training and must use recognised signalling.}
  \bitem{Key takeaway 4}{Ground-bearing capacity assessment must be conducted before positioning mobile cranes, telehandlers, or other heavy plant to prevent overturning or subsidence.}
\end{itemize}

\barquote{In Chapter~\ref{chap:Structural and Temporary Works Risk Assessment}, we turn from this assessment to the next domain in the Prudentia framework.}
